\section{Exercise(s)}
\subsection{Problem A - Alt Tab Command}
In this setting of using our hands to control our computer, the ability to "alt-tab" anytime is a very useful feature that allows the user to navigate the entire workspace in an instant. For this exercise, we want to map this command to a finger or a gesture, so that we can still use the index finger to control the mouse cursor around the screen.

\subsubsection{Resolution A}
Using the Ring Finger (Landmark 16), we can create a simple geometric rule to trigger the Alt-Tab keyboard macro using PyAutoGUI.

\begin{lstlisting}[caption={Exercise A - Alt Tab Implementation}, language=Python, captionpos=b]
rf_point = hand_landmarks.landmark[16] # Ring Finger Tip
rf_base = hand_landmarks.landmark[14]  # Ring Finger PIP Joint (Base)

# If Point is Higher than Base, Ring Finger is raised: Activate Command
if rf_point.y < rf_base.y:             
    # Flag controls Alt key press status, so it doesnt conflict with normal computer use
    if not f_alttab:              
        pyautogui.keyDown('alt')       # Presses Alt key, keeps it pressed
        pyautogui.press('tab')         # Presses Tab key once to enter menu
        f_alttab = True                # Activates Flag
else:
    if f_alttab:
        pyautogui.keyUp('alt')         # Releases Alt key if it was pressed
        f_alttab = False               # Deactivates Flag
\end{lstlisting}

\subsection{Problem B - Screenshot Automation}
Enhance the existing real-time hand-tracking and mouse-control software to trigger a screenshot when a specific gesture is detected. The system needs to recognize when the user fully extends their pinky finger and then execute the \texttt{Windows + Shift + S} keyboard shortcut. Importantly, the program must continue running right after pressing the shortcut. This allows the user to use the existing hand gestures to move the mouse and select the area for the screenshot.

\subsubsection{Resolution B}
To add this feature, we use the \texttt{pyautogui} library to simulate keyboard presses. The logic depends on checking the position of the pinky finger using the MediaPipe hand tracking model. 

According to the model, the tip of the pinky finger is point (landmark) 20. We can figure out if the finger is raised by comparing the Y-coordinate of the tip (landmark 20) with a lower joint, like landmark 18. In images, the Y-axis goes from top to bottom, meaning a smaller Y-value is physically higher on the screen. So, if the tip's Y-value is smaller than the joint's Y-value, the finger is raised.

A major challenge is preventing the shortcut from triggering multiple times a second without pausing the video feed. If we use a simple \texttt{time.sleep()} function, the whole program freezes. This would stop the camera and the mouse control, making it impossible to crop the screenshot. Instead, we use a non-blocking timer with \texttt{time.time()}. This method checks the current time and ensures the shortcut is only pressed once every few seconds. Because the main loop never stops, the user can immediately close their hand (which our Custom AI model predicts as state \texttt{0}) to click and drag the mouse to capture the screen.

The code modification required to implement this logic within the existing software is shown below:
\begin{lstlisting}[caption={Exercise B - Screenshot Automation Implementation}, language=Python, captionpos=b]
import pyautogui
import time

# --- Setup Configuration (Place before the main while loop) ---
pyautogui.PAUSE = 0 
last_action_time = 0
COOLDOWN_SECONDS = 2.0

# ... [Inside the 'for hand_landmarks in results.multi_hand_landmarks:' loop] ...

# 1. Check if the pinky finger is raised
pinky_tip = hand_landmarks.landmark[20].y
pinky_pip = hand_landmarks.landmark[18].y
is_pinky_raised = pinky_tip < pinky_pip

current_time = time.time()

# 2. Trigger the shortcut with a non-blocking timer
if is_pinky_raised and (current_time - last_action_time > COOLDOWN_SECONDS):
    pyautogui.hotkey('win', 'shift', 's')
    last_action_time = current_time # Reset the cooldown timer

# ... [The rest of the existing gesture control logic remains unchanged] ...
\end{lstlisting}

This approach allows the software to take a screenshot and immediately transition back to gesture control, fulfilling the exercise requirements while keeping the application smooth and responsive[cite: 430, 431].